% ============================================================
%  SECTION 3: BIOLOGY OF HONEY RIPENING
% ============================================================

\section{Biology of Honey Ripening: A Continuous Process, Not a Binary Event}

Honey ripening is a complex biotransformation process in which
nectar or honeydew undergoes a series of enzymatic, physicochemical,
and microbiological transformations before reaching a state of
qualitative stability. The widespread equating of honey maturity
with the fact of comb cell capping with wax -- common in beekeeping
practice and veterinary control -- constitutes a simplification that
is inadequate from a biochemical perspective. Capping is one of the
final, morphological indicators of ripening, not synonymous with the
achievement of full chemical and microbiological
stability~\cite{zhang2021,sun2021}.

\subsection{Enzymatic transformation of nectar}

The honey production process is initiated at the very moment nectar
or honeydew is collected by forager bees. Nectar, dominated by
sucrose and containing a high water content, undergoes the first
enzymatic transformation in the honey stomach crop (\textit{honey
stomach crop}) of the insect, where it mixes with secretions
from the bee's pharyngeal
glands~\cite{biochemreactions2022}. The key enzymes initiating
the transformation are:

\begin{itemize}
    \item \textbf{Invertase} ($\beta$-D-fructosidase), catalyzes
          the hydrolysis of sucrose into fructose and glucose;
          produced by the hypopharyngeal glands of worker
          bees~\cite{biochemreactions2022};
          its activity increases gradually from low values
          in immature honey to maximum values in mature honey,
          making it one of the most sensitive indicators of
          the degree of maturity~\cite{zhang2021};
    \item \textbf{Glucose oxidase}, oxidizes glucose to gluconic
          acid with the release of hydrogen peroxide (H$_2$O$_2$);
          gluconic acid is responsible for lowering honey pH to the
          range of 3.5--4.5 and constitutes the dominant organic
          acid in mature honey; H$_2$O$_2$ serves as a natural
          antibacterial agent~\cite{biochemreactions2022,zhang2021};
    \item \textbf{Diastase} ($\alpha$- and $\beta$-amylase),
          catalyzes the hydrolysis of starch and nectar
          polysaccharides into glucose and maltose; diastase
          activity is a parameter standardized by Codex Alimentarius
          ($\geq 8$~Schade units) and serves as an indirect indicator
          of thermal integrity and honey
          maturity~\cite{codex_stan12,biochemreactions2022,zhang2021};
    \item \textbf{Catalase}, decomposes excess H$_2$O$_2$,
          regulating the balance between antibacterial activity
          and protection of the remaining honey components
          from oxidation~\cite{biochemreactions2022}.
\end{itemize}

Studies of enzymatic activity at five successive stages of
transformation -- from floral nectar, through the honey crop of
foragers and house bees, to uncapped and capped hive cells -- showed
a gradual increase in invertase, amylase, and glucose oxidase
activity at each successive stage of
processing~\cite{enzymesrole2023}. These data confirm that
enzymatic transformation proceeds as a \emph{continuum} along
the entire production chain, not as a step-change event occurring
at the moment of cell capping~\cite{enzymesrole2023,zhang2021}.

In parallel with enzymatic transformations, house bees actively
dehydrate honey through hive ventilation: wing beats create a
constant airflow over open comb cells, accelerating water
evaporation and reducing the moisture of the raw material from
values characteristic of nectar to a final approximately 18.5\%
in honey before capping~\cite{zhang2021}. Full maturation under
natural conditions typically takes several days to a few weeks,
with this time potentially being significantly extended under
conditions of high relative
humidity~\cite{zhang2021,tadele2025}.


\subsection{Ripening as a process: metabolomic evidence}

The thesis of the binary nature of honey maturity -- according to
which honey is either immature (uncapped) or mature (capped) --
finds increasingly weak support in current metabolomic research.
Sun et al.~\cite{sun2021} in a study of rapeseed honey
(\textit{Brassica napus}) demonstrated that maturity indicators,
including water content, fructose-to-glucose ratio, invertase and
diastase activity, and organic acid content, change \emph{gradually}
during the ripening process and reach stable values independently
of the moment of cell capping by bees.

Guo et al.~\cite{guo2020} conducted a comprehensive comparative
analysis of the chemical profile and biological activity of mature
and immature honeys using HPLC/QTOF/MS-based metabolomics. The
authors identified significant differences in polyphenol content,
antioxidant activity (DPPH, FRAP), and amino acid profile between
honey collected in an immature state and mature honey; however,
these differences were closely correlated with chemical parameters
(water content, enzymatic activity) rather than solely with
the fact of capping~\cite{guo2020}. Importantly, some samples of
honey collected before full capping exhibited a metabolomic profile
similar to mature honey, provided they met the criterion of low
moisture, which constitutes direct confirmation of the concept
of functional maturity.

Sun et al.~\cite{sun2021} using GC-MS and LC-MS identified
decenedioic acid, a fatty acid of bee origin, as a potential
molecular biomarker differentiating mature from immature honey.
This biomarker, validated for rapeseed, acacia, and jujube honeys,
was significantly more abundant in mature honey ($p < 0{,}01$),
confirming that ripening is a process of gradual accumulation of
specific metabolites, not a binary transition~\cite{sun2021}.

It is worth emphasizing that exposure of honey to high temperatures,
even after full capping, can cause enzyme degradation and effectively
revert quality parameters to values characteristic of immature
honey~\cite{biochemreactions2022}.
Capped honey subjected to inappropriate heat treatment may exhibit
diastase activity below the Codex threshold ($< 8$~Schade units)
and elevated HMF content, demonstrating that the morphological
criterion of capping does not guarantee the maintenance of
functional maturity throughout the supply chain.

This is also confirmed by beekeeping practice, which treats the
fact of cell capping merely as an indicator of potential honey
maturity. On one hand, it sometimes happens that under certain
conditions bees cap comb cells in which honey is not yet mature
(water content exceeds 20\% and shortly after extraction the honey
ferments). On the other hand, common practice involves harvesting
honey from partially capped frames (usually with a threshold of
three-quarters of cells capped in the frame), provided that in
the uncapped cells the honey is mature. The actual indicator of
honey maturity is the water content, which the beekeeper determines
either by refractometry or by other tests (e.g., shaking the frame),
whose result is directly dependent on water content.

\subsection{The concept of functional maturity}
Based on the above review, we propose an operational definition of
\textbf{functional honey maturity} (\textit{functional honey
maturity}), understood as an integrated product state characterized
by the simultaneous fulfillment of three groups of criteria:

\begin{enumerate}
    \item \textbf{Microbiological stability:} absence of active
          fermentation; water activity $a_w < 0{,}60$; moisture
          $\leq 20\%$; osmophilic yeast count below the growth
          threshold under given conditions of temperature
          and storage~\cite{zamora2006}.

    \item \textbf{Chemical stability:} fructose-to-glucose ratio
          (F/G) $\geq 0{,}9$; HMF content $\leq 40$~mg/kg;
          diastase activity $\geq 8$~Schade units; invertase
          activity $\geq 4$~units; gluconic acid content and
          organic acid profile within the reference range
          for the given botanical type~\cite{codex_stan12,zhang2021}.

    \item \textbf{Biological stability:} antioxidant activity
          (DPPH, FRAP) at the level appropriate for the given
          botanical and geographical type; polyphenol and flavonoid
          content not reduced relative to reference values;
          antibacterial activity (assessed by H$_2$O$_2$ production
          and the presence of defensins)
          unchanged~\cite{guo2020}.
\end{enumerate}

This definition deliberately separates the morphological criterion
(cell capping) from the quality criterion, recognizing that honey
may achieve functional maturity before full comb capping (e.g., at
18\% moisture in a temperate climate). This concept is consistent
with the Codex regulatory approach, which, as demonstrated in
Section~2, penalizes the state of fermentation, not
morphology~\cite{codex_stan12,sun2021}.
